\newpage
\section{Contexte et motivations}


Ce projet répond à une double exigence : d’une part, mobiliser un ensemble de compétences techniques clés (programmation orientée objet en Java, structures de données, conception d'interfaces graphiques, planification de trajectoire, etc.) ; d’autre part, s’initier aux méthodes et démarches de génie logicielle à travers la définition de besoins, l’analyse de cas d’usage, la modélisation UML, la répartition modulaire du code et la phase de validation par tests.

La simulation met en scène un environnement semi-réel inspiré à la fois de scénarios de surveillance automatisée et de mécanismes typiques des jeux de stratégie. Les entités principales, à savoir les gardiens et les intrus, sont représentées sous forme d’agents autonomes, capables d’observer, d’analyser et d’agir en fonction de leur environnement immédiat. Les interactions entre agents sont encadrées par un ensemble de règles et d’algorithmes (parcours de graphe, champ de vision, etc.), permettant de simuler une poursuite crédible et dynamique. L’environnement est également constitué d’obstacles fixes et de zones inaccessibles qui renforcent la complexité de la détection et la nécessité d’une planification efficace.

Le choix de ce thème s’est imposé pour plusieurs raisons essentielles. Premièrement, il représente un terrain d’expérimentation stimulant sur les plans technique, algorithmique et créatif. Le joueur ou l’observateur peut immédiatement visualiser les effets de chaque décision logique ou chaque stratégie implémentée, ce qui facilite la compréhension des erreurs et améliore l’apprentissage. Deuxièmement, ce projet incarne un excellent point de convergence entre la théorie et la pratique : des algorithmes stratégique  prennent ici tout leur sens, de gestion d’événements, ou encore d’architecture modulaire.

En outre, la réalisation de ce jeu a permis de nous confronter à des situations réalistes, comme la gestion de threads java , la réalisations de ihm graphique é, ou encore la visualisation fluide d’états changeants dans une interface graphique. Cette confrontation a renforcé notre capacité à résoudre des problèmes complexes de manière structurée et itérative.

Par ailleurs, l’aspect ludique du projet constitue une source de motivation non négligeable : la simulation interactive, qui permet de voir en temps réel les déplacements, la détection et la capture des intrus, donne vie aux concepts abstraits de l’informatique théorique. Ce caractère visuel et dynamique rend les tests plus engageants et suscite un intérêt particulier pour l’optimisation des comportements des agents.

Enfin, au-delà de l’exercice académique, ce projet ouvre une réflexion plus large sur les applications réelles de ce type de simulation. De nombreux domaines professionnels , — font appel à des modèles similaires d’agents intelligents opérant dans des environnements complexes. Il s’agit donc d’une initiation concrète à des problématiques actuelles, qui pourra servir de base à des projets de plus grande envergure dans un futur proche.

Ce contexte, à la fois pédagogique et applicatif, nous a offert l’opportunité de mettre en œuvre une démarche rigoureuse, de structurer une application logicielle complète, et d’adopter une approche de développement incrémental conforme aux standards professionnels. La motivation intrinsèque des membres de l’équipe, combinée à la richesse des sujets abordés, a fait de ce projet une expérience formatrice majeure dans notre parcours universitaire.